\documentclass[a4paper,11pt]{article}
\usepackage{iftex}
\ifPDFTeX
\usepackage[T2A]{fontenc}
\usepackage[utf8]{inputenc}
\usepackage[russian]{babel}
\renewcommand{\familydefault}{\sfdefault}
\else
\usepackage{fontspec}
\setmainfont{Noto Sans}
\usepackage[russian]{babel}
\fi
\usepackage{amsmath,amssymb,mathtools}
\usepackage{icomma}
\usepackage[left=23mm,right=23mm,top=22mm,bottom=23mm]{geometry}
\usepackage{microtype}
\usepackage{needspace}
\usepackage{tikz}
\usetikzlibrary{arrows.meta,positioning,shapes.geometric,calc}
\usepackage{tkz-euclide}
\usepackage{tabularx,booktabs}
\usepackage{enumitem}
\usepackage{fancyhdr}
\usepackage{setspace}
\usepackage{xcolor}
\usepackage{hyperref}
\definecolor{accent}{HTML}{008080}
\definecolor{darkaccent}{HTML}{004D4D}
\definecolor{textgray}{HTML}{333333}
\definecolor{softgray}{HTML}{6C7478}
\hypersetup{colorlinks=true,linkcolor=darkaccent,urlcolor=accent}
\setstretch{1.15}
\setlength{\headheight}{22pt}
\setlength{\parindent}{0pt}
\setlength{\parskip}{5pt}
\setlist[itemize]{leftmargin=6mm,itemsep=1.5pt,topsep=3pt}
\setlist[enumerate]{leftmargin=7mm,itemsep=7pt,topsep=4pt}
\pagestyle{fancy}
\fancyhf{}
\fancyhead[L]{\small\color{softgray} Геометрия и линейная функция}
\fancyhead[R]{\small\color{softgray} Кирилл Зиновьев}
\fancyfoot[C]{\small\thepage}
\renewcommand{\headrulewidth}{0.3pt}
\renewcommand{\headrule}{\hbox to\headwidth{\color{accent}\leaders\hrule height \headrulewidth\hfill}}
\newcommand{\secname}[1]{\vspace{6pt}{\Large\color{darkaccent}#1\par}\vspace{3pt}\noindent{\color{accent}\rule{\linewidth}{0.45pt}}\vspace{1pt}}
\newcommand{\subname}[1]{{\color{darkaccent}\bfseries #1}\par}
\newcommand{\idea}[1]{\textcolor{darkaccent}{\textbf{Главная идея.}} #1}
\newcommand{\warn}[1]{\textcolor{darkaccent}{\textbf{Обратите внимание.}} #1}
\newcommand{\answerrow}[2]{#1 & #2\\}
\begin{document}
\begin{titlepage}
\thispagestyle{empty}
\centering
\vspace*{37mm}
{\large\color{darkaccent} МАТЕМАТИКА\par}
\vspace{22mm}
{\huge Чек-лист\par}
\vspace{13mm}
{\LARGE\color{accent} Углы и равенство треугольников\par}
\vspace{3mm}
{\LARGE\color{accent} Линейная функция и точки графика\par}
\vspace{20mm}
{\normalsize Повторение материала занятия · 7 класс\par}
\vfill
{\large Кирилл Зиновьев\par}
\vspace{12mm}
{\normalsize 19 ноября 2026 г.\par}
\vspace{18mm}
{\small Лёвин Артём Александрович\par}
{\small эксклюзивно для Кирилла Зиновьева\par}
\vspace{15mm}
\begin{tikzpicture}[x=1cm,y=1cm]
\draw[accent, line width=0.8pt] (0,0) .. controls (2.1,0.65) and (3.5,-0.55) .. (5.0,0.02)
.. controls (6.8,0.65) and (8.4,-0.40) .. (10.1,0.15);
\end{tikzpicture}
\vspace{5mm}
\end{titlepage}

\secname{1. Как правильно обозначать углы}
В задачах с несколькими треугольниками особенно важно точно указывать, о каком угле идёт речь.

\begin{itemize}
\item Угол \(\angle ADB\) имеет вершину \(D\), его стороны --- лучи \(DA\) и \(DB\).
\item Угол \(\angle ADC\) тоже имеет вершину \(D\), но его стороны --- лучи \(DA\) и \(DC\).
\item Поэтому \(\angle ADB\) и \(\angle ADC\) в общем случае \emph{разные углы}. Буква вершины всегда стоит посередине.
\end{itemize}
\warn{У треугольников может быть общая сторона; «общая буква» ещё не означает общий угол.}

\secname{2. Смежные, вертикальные и углы треугольника}
\begin{minipage}[t]{0.57\linewidth}
\vspace{0pt}
\subname{Смежные углы}
Имеют общую сторону, а две другие стороны образуют прямую:
\[
\alpha+\beta=180^\circ.
\]
\subname{Вертикальные углы}
Образуются при пересечении двух прямых и равны:
\[
\angle AOB=\angle COD.
\]
\subname{Сумма углов треугольника}
\[
\angle A+\angle B+\angle C=180^\circ.
\]
\end{minipage}\hfill
\begin{minipage}[t]{0.39\linewidth}
\vspace{0pt}\centering
\begin{tikzpicture}[scale=.76]
\tkzDefPoint(0,0){O}
\tkzDefPoint(2.9,1.6){A}
\tkzDefPoint(-2.9,-1.6){C}
\tkzDefPoint(-2.9,1.6){B}
\tkzDefPoint(2.9,-1.6){D}
\tkzDrawSegments(A,C B,D)
\tkzDrawPoints(O)
\tkzLabelPoints[above](A,B)
\tkzLabelPoints[below](C,D)
\tkzLabelPoints[below right](O)
\end{tikzpicture}
\vspace{5pt}
{\small При пересечении прямых вертикальные углы равны.}
\end{minipage}

\subname{Пример: два свойства подряд}
Пусть \(\angle AOC=100^\circ\), причём \(OB\) и \(OC\) --- противоположные лучи, а \(OA\) и \(OD\) --- противоположные лучи. Тогда
\[
\angle AOB=180^\circ-100^\circ=80^\circ,
\qquad \angle COD=80^\circ.
\]
Последнее равенство следует из свойства вертикальных углов.

\Needspace{.36\textheight}
\secname{3. Равнобедренный треугольник}
Если \(AB=BC\), то \(AC\) --- основание треугольника \(ABC\), а углы при основании равны:
\[
\angle BAC=\angle BCA.
\]
\begin{minipage}[t]{0.62\linewidth}
\vspace{0pt}
\subname{Пример}
В равнобедренном треугольнике с основанием \(AC\) угол при вершине \(B\) равен \(42^\circ\). Тогда
\[
\angle A=\angle C=\frac{180^\circ-42^\circ}{2}
=69^\circ.
\]
Делим на два, поскольку оставшиеся углы равны.
\end{minipage}\hfill
\begin{minipage}[t]{0.34\linewidth}
\vspace{0pt}\centering
\begin{tikzpicture}[scale=.79]
\tkzDefPoint(-2.3,0){A}
\tkzDefPoint(2.3,0){C}
\tkzDefPoint(0,3.2){B}
\tkzDrawPolygon(A,B,C)
\tkzMarkSegments[mark=|,size=3pt](A,B B,C)
\tkzLabelPoints[below](A,C)
\tkzLabelPoints[above](B)
\end{tikzpicture}
\end{minipage}

\vspace{5mm}
\secname{4. Второй признак равенства треугольников}
\idea{Если сторона одного треугольника и два прилежащих к ней угла соответственно равны стороне и двум прилежащим к ней углам другого треугольника, такие треугольники равны.}

\subname{Пример, близкий к задаче занятия}
Пусть прямые \(AD\) и \(BC\) пересекаются в \(O\), \(BO=DO\), \(\angle ABC=45^\circ\), \(\angle BCD=55^\circ\), \(\angle AOC=100^\circ\). Найдём \(\angle CDO\) и докажем равенство треугольников \(ABO\) и \(CDO\).

\begin{center}
\begin{tikzpicture}[scale=.70]
\tkzDefPoint(0,0){O}
\tkzDefPoint(1.665,1.984){A}
\tkzDefPoint(-1.928,2.298){B}
\tkzDefPoint(1.665,-1.984){C}
\tkzDefPoint(-1.928,-2.298){D}
\tkzDrawSegments(A,B B,O O,A C,D D,O O,C)
\tkzMarkSegments[mark=|,size=4pt](B,O D,O)
\tkzDrawPoints(O)
\tkzLabelPoints[above](A,B)
\tkzLabelPoints[below](C,D)
\tkzLabelPoints[right](O)
\node[font=\small] at (-2.48,1.80){\(45^\circ\)};
\node[font=\small] at (2.31,-1.83){\(55^\circ\)};
\node[font=\small] at (1.28,0){\(100^\circ\)};
\end{tikzpicture}
\end{center}
\textbf{Шаг 1.} \(\angle AOB=180^\circ-100^\circ=80^\circ\) (смежные углы).

\textbf{Шаг 2.} \(\angle COD=80^\circ\) (вертикальные углы).

\textbf{Шаг 3.} Из треугольника \(CDO\):
\[
\angle CDO=180^\circ-80^\circ-55^\circ=45^\circ.
\]
\textbf{Шаг 4.} \(\angle ABO=45^\circ=\angle CDO\), \(\angle AOB=\angle COD=80^\circ\), \(BO=DO\). Значит,
\[
\boxed{\triangle ABO\cong\triangle CDO}
\]
по второму признаку: равны сторона и два прилежащих к ней угла.

\warn{При доказательстве всегда называйте соответствующие вершины в согласованном порядке.}

\secname{5. Как находить углы в прямоугольных треугольниках}
Угол прямоугольного треугольника, отличный от прямого, находится вычитанием из \(90^\circ\):
\[
\angle BAD=90^\circ-\angle ADB.
\]
При сравнении двух треугольников, имеющих общую сторону \(AD\), сначала вычислите прилежащие к этой стороне углы, затем проверьте второй признак.

\vspace{7mm}
\secname{6. Линейная функция и координатные оси}
Линейная функция задаётся формулой
\[
\boxed{y=kx+b},\qquad k,b\in\mathbb{R}.
\]
Графиком служит прямая (при \(k=0\) --- горизонтальная). Если требуется найти пересечения с осями, строить график необязательно.

\renewcommand{\arraystretch}{1.3}
\begin{center}
\begin{tabularx}{\linewidth}{@{}p{.28\linewidth}XX@{}}
\toprule
\textbf{Что ищем} & \textbf{Что известно} & \textbf{Что делаем}\\
\midrule
Пересечение с \(Ox\) & \(y=0\) & Подставляем \(0\) вместо \(y\), решаем уравнение для \(x\). Точка \((x;0)\).\\
Пересечение с \(Oy\) & \(x=0\) & Подставляем \(0\) вместо \(x\), находим \(y\). Точка \((0;y)\).\\
\bottomrule
\end{tabularx}
\end{center}

\subname{Пример: обе точки пересечения}
Для \(y=-2{,}4x+9{,}6\):
\begin{align*}
Ox:\quad&0=-2{,}4x+9{,}6
\quad\Rightarrow\quad x=4
\quad\Rightarrow\quad A(4;0);\\
Oy:\quad&y=-2{,}4\cdot0+9{,}6=9{,}6
\quad\Rightarrow\quad B(0;9{,}6).
\end{align*}
\warn{Ноль определяется выбранной осью. При поиске пересечения с \(Ox\) обнуляется \(y\), с \(Oy\) --- \(x\).}

\subname{Различайте три формулировки}
\begin{itemize}
\item \textbf{«В какой точке график пересекает ось \(Ox\)?»} --- подставить \(y=0\).
\item \textbf{«В какой точке график пересекает ось \(Oy\)?»} --- подставить \(x=0\).
\item \textbf{«В какой точке пересекаются графики двух функций?»} --- приравнять выражения для \(y\) и решить полученное уравнение. Это другой тип задачи.
\end{itemize}

\Needspace{.35\textheight}
\secname{7. Проходит ли график через заданную точку?}
Точка \(M(x_0;y_0)\) лежит на графике \(y=kx+b\) тогда и только тогда, когда
\[
\boxed{y_0=kx_0+b.}
\]
\textbf{Пример.} Для \(y=1{,}2x-7\) и точки \(M(100;113)\):
\[
1{,}2\cdot100-7=120-7=113.
\]
Вычисленное значение совпало с заданной второй координатой. Значит, точка \(M\) лежит на графике.

\warn{Подставьте \(x\) точки в формулу функции и сравните полученное значение с \(y\) той же точки.}

\secname{8. Типичные ошибки и контроль понимания}
\begin{itemize}
\item \(\square\) В записи \(\angle ADB\) я называю вершиной \(D\) и различаю углы \(ADB\) и \(ADC\).
\item \(\square\) Для смежных углов использую сумму \(180^\circ\), а для вертикальных --- равенство.
\item \(\square\) В равнобедренном треугольнике делю оставшуюся сумму углов на два, потому что углы при основании равны.
\item \(\square\) При втором признаке проверяю именно сторону и \emph{два прилежащих к ней угла}.
\item \(\square\) Для пересечения с \(Ox\) полагаю \(y=0\), для \(Oy\) --- \(x=0\).
\item \(\square\) Координаты записываю в порядке \((x;y)\) и подписываю ось.
\item \(\square\) Проверяя точку, сравниваю вычисленное \(y\) с её второй координатой.
\end{itemize}

\subname{Как выбрать метод по условию}
\begin{center}
\begin{tabularx}{\linewidth}{@{}p{.44\linewidth}X@{}}
\toprule
\textbf{Сигнал в условии} & \textbf{Первое действие}\\
\midrule
Даны пересекающиеся прямые & Найдите вертикальные и смежные углы.\\
Дан равнобедренный треугольник & Отметьте равные стороны и углы при основании.\\
Нужно доказать равенство треугольников & Установите соответствие вершин и проверьте признак.\\
Нужна точка пересечения с осью & Подставьте ноль в нужную координату.\\
Нужно проверить точку графика & Подставьте её абсциссу в формулу и сравните ординату.\\
\bottomrule
\end{tabularx}
\end{center}

\clearpage
\thispagestyle{fancy}
\begin{center}
{\Large\color{darkaccent} Алгоритм: принадлежит ли точка графику?}\par
\vspace{14mm}
\begin{tikzpicture}[
font=\small,
node distance=12mm,
start/.style={ellipse,draw=accent,thick,minimum width=45mm,minimum height=11mm,align=center},
act/.style={rectangle,rounded corners=2pt,draw=accent,thick,text width=91mm,minimum height=13mm,align=center,inner sep=4pt},
dec/.style={diamond,aspect=2.6,draw=darkaccent,thick,align=center,text width=37mm,inner sep=1pt},
last/.style={rectangle,rounded corners=2pt,draw=accent,thick,text width=46mm,minimum height=20mm,align=center,inner sep=4pt},
arr/.style={-{Stealth[length=3mm]},thick,color=darkaccent}
]
\node[start] (s) {Дана точка $M(x_0;y_0)$};
\node[act,below=of s] (f) {Запишите формулу функции $y=kx+b$};
\node[act,below=of f] (calc) {Подставьте $x_0$ вместо $x$ и вычислите $y_{\text{расч}}=kx_0+b$};
\node[dec,below=18mm of calc] (q) {$y_{\text{расч}}=y_0$?};
\node[last,below left=24mm and 17mm of q] (yes) {Да. Точка принадлежит графику};
\node[last,below right=24mm and 17mm of q] (no) {Нет. Точка не принадлежит графику};
\draw[arr] (s)--(f);
\draw[arr] (f)--(calc);
\draw[arr] (calc)--(q);
\draw[arr] (q.south west) -| node[pos=.20,above left]{да} (yes.north);
\draw[arr] (q.south east) -| node[pos=.20,above right]{нет} (no.north);
\end{tikzpicture}
\end{center}
\clearpage

\clearpage
\secname{Тренировочный блок}
\textcolor{softgray}{Путь А · 10 заданий · от базовых приёмов к сочетанию нескольких идей.}

\subname{Геометрия}
\begin{enumerate}
\item Углы \(AOB\) и \(BOC\) смежные. \(\angle AOB=116^\circ\). Найдите \(\angle BOC\).
\item Прямые \(AC\) и \(BD\) пересекаются в точке \(O\). Известно, что \(\angle AOB=67^\circ\). Найдите \(\angle COD\) и \(\angle BOC\).
\item В равнобедренном треугольнике \(ABC\) \(AB=BC\), а \(\angle ABC=48^\circ\). Найдите два других угла.
\item В равнобедренном треугольнике \(ABC\) с основанием \(AC\) дано \(\angle BAC=72^\circ\). Найдите \(\angle ABC\).
\item Треугольники \(ABD\) и \(ACD\) прямоугольные: \(\angle ABD=\angle ACD=90^\circ\). Они имеют общую гипотенузу \(AD\). Дано \(\angle ADB=35^\circ\), \(\angle ADC=55^\circ\). Вычислите \(\angle BAD\) и \(\angle CAD\). Докажите, что \(\triangle ABD\cong\triangle DCA\), используя второй признак.
\end{enumerate}

\subname{Алгебра}
\begin{enumerate}[start=6]
\item Без построения найдите точки пересечения графика \(y=-3x+12\) с осями \(Ox\) и \(Oy\).
\item Без построения найдите точки пересечения графика \(y=0{,}5x-4\) с осями \(Ox\) и \(Oy\).
\item Без построения найдите точки пересечения графика \(y=-1{,}25x+5\) с осями \(Ox\) и \(Oy\).
\item Определите, какие точки лежат на графике функции \(y=1{,}2x-7\): \(P(100;113)\), \(Q(10;6)\). Обоснуйте ответ подстановкой.
\item Для функции \(y=0{,}4x-12\) найдите точки пересечения с осями; затем определите, лежат ли на графике точки \(M(5;-10)\) и \(N(-5;-13)\).
\end{enumerate}

\clearpage
\secname{Ответы к тренировочному блоку}
\small
\renewcommand{\arraystretch}{1.36}
\begin{center}
\begin{tabularx}{\linewidth}{@{}p{.10\linewidth}X@{}}
\toprule
\textbf{№} & \textbf{Ответ / результат}\\
\midrule
1 & \(64^\circ\).\\
2 & \(\angle COD=67^\circ\), \(\angle BOC=113^\circ\).\\
3 & \(\angle BAC=\angle BCA=66^\circ\).\\
4 & \(\angle ABC=36^\circ\).\\
5 & \(\angle BAD=55^\circ\), \(\angle CAD=35^\circ\). Углы при стороне \(AD\) попарно равны; \(\triangle ABD\cong\triangle DCA\) по второму признаку.\\
6 & \(Ox:\ (4;0)\); \(Oy:\ (0;12)\).\\
7 & \(Ox:\ (8;0)\); \(Oy:\ (0;-4)\).\\
8 & \(Ox:\ (4;0)\); \(Oy:\ (0;5)\).\\
9 & \(P\) лежит на графике, \(Q\) не лежит: при \(x=10\) получается \(y=5\).\\
10 & \(Ox:\ (30;0)\); \(Oy:\ (0;-12)\). \(M\) лежит на графике, \(N\) не лежит: при \(x=-5\) получается \(y=-14\).\\
\bottomrule
\end{tabularx}
\end{center}
\normalsize
\vfill
\begin{center}
{\color{darkaccent}\rule{.68\linewidth}{.45pt}}\\[7pt]
{\small Лёвин Артём Александрович\par}
{\small эксклюзивно для Кирилла Зиновьева\par}
\end{center}
\end{document}
